\documentclass[12pt]{article}

%% Basic document formatting
\usepackage{amsmath, amsthm, amssymb, amsfonts}
\usepackage{mathtools}
\usepackage{xspace}
\usepackage{thmtools}
\usepackage{graphicx}
\usepackage{setspace}
\usepackage{fancyhdr}
\usepackage{titling}
\usepackage[left=0.4in,right=0.4in,top=1in,bottom=1in]{geometry}
\usepackage{float}
\usepackage{hyperref}
\usepackage{mathrsfs}
\usepackage{tabularx}
\usepackage[utf8]{inputenc}
\usepackage[english]{babel}
\usepackage{framed}
\usepackage[dvipsnames]{xcolor}
\usepackage{environ}
\usepackage{tcolorbox}
\tcbuselibrary{theorems,skins,breakable}

\usepackage{enumitem}
\setlist[enumerate]{leftmargin=*}
\setlist[enumerate,1]{labelindent=\parindent}
\setlist[enumerate,2]{labelindent=0pt}

% Blackboard Bold
\newcommand{\Z}{\mathbb{Z}}                 % integers
\newcommand{\Zpl}{\mathbb{Z}_{+}}           % positive integers
\newcommand{\Znn}{\mathbb{Z}_{\geq 0}}      % nonnegative integers. 
\newcommand{\Q}{\mathbb{Q}}                 % rationals
\newcommand{\Qpl}{\mathbb{Q}_{+}}           % positive Rationals
\newcommand{\R}{\mathbb{R}}                 % reals
\newcommand{\Rpl}{\mathbb{R}_{+}}           % positive Reals
\newcommand{\C}{\mathbb{C}}                 % complex numbers
\newcommand{\F}{\mathbb{F}}                 % field

% Words
\newcommand{\st}{\text{ such that }}
\newcommand{\wrt}{\text{ with respect to }}
\newcommand{\with}{\text{ with }}
\newcommand{\ie}{\text{, i.e. }}
\newcommand*{\ditto}{\text{---}\texttt{"}\text{---}}

% Operators
\newcommand{\abs}[1]{\left|#1\right|}                   % absolute value
\newcommand{\norm}[1]{\abs{\abs{#1}}}
\newcommand{\floor}[1]{\left\lfloor #1 \right\rfloor}   % floor
\newcommand{\ceil}[1]{\left\lceil #1 \right\rceil}      % ceiling
\newcommand{\powset}[1]{\mathscr{P}(#1)}

% Category Theory 
\renewcommand{\hom}{\text{Hom}}
\newcommand{\homspace}[3]{\hom_{#1}(#2, #3)}
\newcommand{\coproduct}[9]{
  \begin{tikzcd}[ampersand replacement=\&]
      \& \& #1 \arrow[dll, bend right, "#6"] \arrow[dl, "#5"] \\
   #4 \& #3 \arrow[l, "#9"] \\
      \& \& #2 \arrow[ull, bend left, "#8"] \arrow[ul, "#7"]
  \end{tikzcd}
}

\newcommand{\product}[9]{ 
  \begin{tikzcd}[ampersand replacement=\&]
      \& \& #3 \\
      #1 \arrow[r, "#7"] \arrow[urr, bend left, "#5"] \arrow[drr, bend right, "#6"] \& #2 \arrow[ur, "#8"] \arrow[dr, "#9"] \\ 
      \& \& #4
  \end{tikzcd}
}


% Algebra 
\newcommand{\Syl}[2]{\text{Syl}_{#1}{(#2)}}           % set of Sylow-p subgroups 
\newcommand{\<}{\langle}                % \<x\>, subgroup generated by x 
\renewcommand{\>}{\rangle}
\renewcommand{\index}[2]{[#1 : #2]}
\newcommand{\id}{\text{id}}             % identity element
\newcommand{\lhdneq}{%
  \mathrel{\ooalign{$\lneq$\cr\raise.22ex\hbox{$\lhd$}\cr}}}
\newcommand{\order}[1]{\text{o}(#1)}    % order of an element
\newcommand{\spec}{\operatorname{Spec}}
\newcommand{\rad}{\operatorname{rad}}   % radical of an ideal 
\newcommand{\sym}[1]{\mathfrak{S}_{#1}}
\newcommand{\aut}[1]{\text{Aut}(#1)} 
\newcommand{\gal}[2]{\text{Gal}(#1 / #2)} 
\newcommand{\inn}[1]{\text{Inn}(#1)} 
\let\oldcong\cong
\let\oldequiv\equiv
\renewcommand{\cong}{\oldequiv}
\renewcommand{\equiv}{\oldcong}

\newcommand{\triangleleftneq}{\mathrel{\ooalign{%
  $\trianglelefteq$\cr
  \hidewidth\raise0.06ex\hbox{$\not\mkern4mu$}\hidewidth\cr
}}}

\makeatletter % cycle
\newcommand{\cyc}[1]{(\mathbf{\cyc@process#1\relax})}
\def\cyc@process#1#2\relax{%
	#1%
	\ifx\relax#2\relax
	\else
		\,\cyc@process#2\relax
	\fi
}
\makeatother

\newcommand{\GL}[2]{\text{GL}_{#1}(#2)}                             % general linear group
\newcommand{\SL}[2]{\text{SL}_{#1}(#2)}                             % special linear group
\newcommand{\gl}[2]{\mathfrak{gl}_{#1}(#2)}
\renewcommand{\sl}[2]{\mathfrak{sl}_{#1}(#2)}
\newcommand{\so}[2]{\mathfrak{so}_{#1}(#2)}
\newcommand{\su}[1]{\mathfrak{su}(#1)}

% Linear Algebra
\newcommand{\bmat}[1]{\begin{bmatrix}#1\end{bmatrix}}   % bracketed matrix
\newcommand{\rank}{\operatorname{rank}}                 % rank
\newcommand{\nullity}{\operatorname{nullity}}           % nullity
\newcommand{\tr}{\operatorname{tr}}

% Combinatorics
\newcommand{\qnum}[1]{\left[#1\right]_q}                % q-analog
\newcommand{\qfact}[1]{\left[#1\right]!_q}              % q-analog factorial 
\newcommand{\qbinom}[2]{\binom{#1}{#2}_q}               % Gaussian binomial coefficient

% Topology/Analysis
\newcommand{\ball}[2]{\text{B}_{#1}(#2)}  % B_r(x): open r-balls around x
\newcommand{\diam}{\text{diam}}           % diamter of a set in metric space 
\newcommand{\lowint}[2]{\underline{\int_{#1}^{#2}}}
\newcommand{\upint}[2]{\overline{\int}_{#1}^{#2}}

% Misc Notation
\newcommand{\oneton}[1]{\{1, \ldots, #1\}}
\newcommand{\defeq}{\vcentcolon=}   % :=
\newcommand{\eqdef}{=\vcentcolon}   % =: 
\renewcommand{\bf}[1]{\textbf{#1}}
\newcommand{\im}{\operatorname{im}}
\newcommand{\fnrest}[2]{\left.#1\right|_{#2}}
\newcommand{\isom}{\overset{\sim}{\rightarrow}} 


\renewcommand{\thesection}{\arabic{section}.}
\renewcommand{\thesubsection}{(\alph{subsection})}

\newtheorem{claim}{Claim}
\newtheorem{theorem}{Theorem}
\newtheorem{proposition}{Proposition}
\newtheorem*{lemma}{Lemma}

%% Headers & title setup
\newcommand{\course}{Math140B}
\newcommand{\myname}{Jay Ser}
\setlength{\headheight}{14.5pt}
\pagestyle{fancy}
\fancyhf{}
\renewcommand{\headrulewidth}{0.4pt}
\lhead{\course}
\rhead{\myname}
\cfoot{\thepage}
\setlength{\droptitle}{-4em} 
\title{\course\ - HW \#9}
\author{\myname}
\date{2026.03.08} 

\begin{document}
\maketitle
\thispagestyle{fancy}

%------------------------------------------------------------------------------%

\section*{Q4} 
\subsection{}
Notice 
$$\lim_{x \rightarrow 0} \frac{b^x - 1}{x} = f'(0)$$ 
where $f(x) = b^x = e^{x \log b}$. 
Then 
$$f'(x) = e^{x \log b} \log b = b^x \log b$$
So $f'(0) = \log b$, as was to be shown. 

\subsection{}
By the L'Hospital rule, 
$$\lim_{x \rightarrow 0} \frac{\log(1 + x)}{x} = \lim_{x \rightarrow 0} \frac{1}{1 + x} = 1$$

\subsection{}
Part (c) follows from part (b). 
$$\lim_{x \rightarrow 0} (1 + x)^{1/x} = \lim_{x \rightarrow 0} e^{1/x \log{(1 + x)}} = e$$

\subsection{} 
Proceed as in the proof for Theorem 3.31 of Rudin. 
Fix $x \in \R$ and let 
$$f_n(x) = (1 + \frac{x}{n})^n$$ 
By the Binomial Theorem, 
\begin{align}
  f_n(x) & = \sum_{k = 0}^{n} \binom{n}{k} (\frac{x}{n})^k \nonumber \\ 
         & = \sum_{k = 0}^{n} \frac{x^k}{k!} \frac{n(n - 1) \cdots (n - k + 1)}{n^k} \nonumber \\ 
         & = \sum_{k = 0}^{n} \frac{x^k}{k!} (1 - \frac{1}{n})(1 - \frac{2}{n}) \cdots (1 - \frac{k - 1}{n}) \label{eq:binom} \\ 
         & \leq \sum_{k = 0}^{n} \frac{x^k}{k!} \nonumber
\end{align}
which gives  
$$\limsup_{n \rightarrow \infty} f_n(x) \leq e^x$$ 
Next, if $n \geq m$, then by \eqref{eq:binom},  
$$f_n(x) \geq \sum_{k = 0}^{m} \frac{x^k}{k!} (1 - \frac{1}{n})(1 - \frac{2}{n}) \cdots (1 - \frac{k - 1}{n}) \geq \sum_{k = 0}^{m} \frac{x^k}{k!}$$ 
Fixing $m$ and letting $n \rightarrow \infty$, 
$$\liminf_{n \rightarrow \infty} f_n(x) \geq \sum_{k = 0}^{m} \frac{x^k}{k!}$$
Letting $m \rightarrow \infty$, obtain 
$$\liminf_{n \rightarrow \infty} f_n(x) \geq e^x$$ 
This gives the desired result. 

\section*{Q5}
\setcounter{subsection}{0}
\subsection{}
Recall from Q4 (c) that $\lim_{x \rightarrow 0} \frac{e - (1 + x)^{1/x}}{x}$ and $(1 + x)^{1/x}$ is differentiable everywhere except $x = 0$, so L'Hospital's rule applies:  
\begin{equation} \label{eq:5_res}
  \lim_{x \rightarrow 0} \frac{e - (1 + x)^{1/x}}{x} = - \lim_{x \rightarrow 0} \frac{dy}{dx}
\end{equation}
where $y = (1 + x)^{1/x}$. 
$\log y = -x \log(1 + x)$, so 
$$\frac{1}{y} \frac{dy}{dx} = \frac{d}{dx} \log y = \frac{1}{x(1 + x)} - \frac{\log(1 + x)}{x^2} $$
This gives 
$$\frac{dy}{dx} = -(1 + x)^{1/x} \left(\frac{1}{x(1 + x)} - \frac{\log(1 + x)}{x^2}\right)$$
As $x \rightarrow 0$, the first term approaches $e$. 
L'Hospital's rule can be applied a few more times to calculate limit of the second term as $x \rightarrow 0$: 
\begin{align*}
  \lim_{x \rightarrow 0} \left[ \frac{1}{x(1 + x)} - \frac{\log(1 + x)}{x^2} \right] 
  & = \lim_{x \rightarrow 0} \frac{x - (1 + x)\log(1 + x)}{x^2 (1 + x)} \\ 
  & = \lim_{x \rightarrow 0} \frac{1 - \frac{1 + x}{1 + x} - \log(1 + x)}{2x + 3x^2} \\ 
  & = \lim_{x \rightarrow 0} \frac{-1/(1 + x)}{2 + 6x} = -1/2
\end{align*}
Recapitulating, 
$$-\lim_{x \rightarrow 0} \frac{dy}{dx} = e / 2$$ 

\subsection{}
I will show 
$$\lim_{n \rightarrow \infty} \frac{n}{\log{n}} (n^{1/n} - 1) = \lim_{x \rightarrow \infty} \frac{x}{\log{x}}(x^{1/x} - 1) = 1$$ 
where the first equality is trivial. 
Recall that $\lim_{x \rightarrow \infty} \log{x}/x = 0$. 
Applying L'Hospital's rule, 
\begin{align*}
  \lim_{x \rightarrow \infty} \frac{x}{\log{x}}(x^{1/x} - 1) 
  & = \lim_{x \rightarrow \infty} \frac{e^{\log{x} / x} - 1}{\log{x} / x} \\ 
  & = \lim_{x \rightarrow \infty} \frac{y e^{\log{x} / x}}{y} \\ 
  & = e^0 = 1
\end{align*}
where $y = \frac{d}{dx} \log{x} / x$. 

\section*{Q9}
\setcounter{subsection}{0}
\subsection{}
I show that $s_N - \log{N}$ is a monotonic decreasing and bounded sequence.
To see that the sequence is monotonic decreasing, 
\begin{align*}
  s_{N + 1} - \log{N + 1} - (s_{N} - \log{N}) & = \frac{1}{N + 1} - \int_{N}^{N + 1} \frac{dt}{t} \\ 
                                              & < \frac{1}{N + 1} - \frac{1}{N + 1} \\ 
                                              & = 0
\end{align*}
where the inequality comes from letting $P = \{N, N + 1\}$ be a one-block partition of $[N, N + 1]$ and the fact that $L(P, \frac{1}{x}) = \frac{1}{N + 1}$ since $\frac{1}{x}$ is a decreasing function. 

To show that $s_N - \log{N}$ is bounded, it suffices to show that each term in the sequence is positive. 
Indeed, for any $N$, 
\begin{align*}
  s_N - \log{N} & = s_N - \int_{1}^{N} \frac{dt}{t} \\ 
                & > s_N - \sum_{k = 1}^{N - 1} \frac{1}{k} \\ 
                & = \frac{1}{N} > 0
\end{align*}
where the inequality arises from letting $P = \{1, 2, \ldots, N - 1, N\}$ be a partition of $[1, N]$ and the fact that $U(P, \frac{1}{x}) = \sum_{k = 1}^{N - 1} \frac{1}{k}$ since $\frac{1}{x}$ is a decreasing function. 

I have shown that $s_N - \log{N}$ is monotonic decreasing and bounded. 
Hence, it is a converging sequence. 

\section*{Q11}
Notice 
\begin{equation} \label{eq:one}
  1 = t\int_{0}^{\infty} e^{-tx} dx
\end{equation}
The desired result is equivalent to 
\begin{equation} \label{eq:end_res} 
  \lim_{t \rightarrow 0} \abs{t \int_{0}^{\infty} e^{-tx} f(x) dx - 1} = 0
\end{equation}

Fix $\epsilon > 0$. 
Let $A \in \R$ be the value such that $\abs{f(x) - 1} < \epsilon$ for all $x \geq A$. 
Denote $M \defeq \sup_{x \in [0, A]} \abs{f(x) - 1}$. 
Combining \eqref{eq:one} and \eqref{eq:end_res}, obtain 
\begin{align}
  \limsup_{t \rightarrow 0} \abs{t \int_{0}^{\infty} e^{-tx} dx - 1} 
  & = \limsup_{t \rightarrow 0} \abs{t \int_{0}^{\infty} e^{-tx} (f(x) - 1)dx} \nonumber \\ 
  & \leq \lim_{t \rightarrow 0} \abs{t} \left[ \int_{0}^{A} M e^{-tx} dx + \int_{A}^{\infty} \epsilon e^{-tx}dx \right] \label{eq:pos_t} \\ 
  & = \lim_{t \rightarrow 0} \abs{t} \left[ \left.-\frac{M}{t}e^{-tx} \right|_{0}^{A} + \left. -\frac{\epsilon}{t}e^{-tx}\right|_{A}^{\infty}\right] \nonumber \\ 
  & = \lim_{t \rightarrow 0} \left[M(e^{-Ax} - 1) + \epsilon e^{-At} \right] \label{eq:cancel} \\
  & = \epsilon \nonumber
\end{align}
where in \eqref{eq:cancel}, the $\abs{t}$ and $\frac{1}{t}$ cancel because only the nonnegative values of $t$ are considered when integrating $e^{-tx}$ in \eqref{eq:pos_t}. 
Since $\epsilon$ was chosen to be arbitrary, 
$$\limsup_{t \rightarrow 0} \abs{t \int_{0}^{\infty} e^{-tx} dx - 1} = 0$$
which proves \eqref{eq:end_res}. 

\section*{Q22}
Let 
\begin{equation} \label{eq:22_series}
  f(x) = 1 + \sum_{n = 1}^{\infty} \frac{\alpha (\alpha - 1) \cdots (\alpha - n + 1)}{n!} x^n
\end{equation}
The series in \eqref{eq:22_series} converges with a radius of convergence of 1 by the ratio test: 
$$\limsup_{n \rightarrow \infty} \abs{\frac{x(\alpha - n)}{n + 1}} < 1 \iff \abs{x} < 1$$
To obtain the differential equation, first calculate 
$$f'(x) = \sum_{n = 1}^{\infty} \frac{\alpha(\alpha - 1) \cdots (\alpha - n + 1)}{(n - 1)!} x^{n - 1} = \sum_{n = 0}^{\infty} \frac{\alpha(\alpha - 1) \cdots (\alpha - n)}{n!} x^n$$
Recall the notation $\binom{\alpha}{n} = \frac{\alpha(\alpha - 1) \cdots (\alpha - n + 1)}{n!}$.
A result from combinatorics is that $\binom{\alpha}{n} = \binom{\alpha - 1}{n - 1} + \binom{\alpha - 1}{n}$. 
Then 
\begin{align*}
  \frac{1 + x}{\alpha} f'(x) & = \sum_{n = 1}^{\infty} \frac{(\alpha - 1)(\alpha - 2) \cdots (\alpha - n + 1)}{(n - 1)!} x^n + \sum_{n = 0}^{\infty} \frac{(\alpha - 1) (\alpha - 2) \cdots (\alpha - n)}{n!} x^n  \\ 
                             & = 1 + \sum_{n = 1}^{\infty} (\binom{\alpha - 1}{n - 1} + \binom{\alpha - 1}{n}) x^n \\ 
                             & = 1 + \sum_{n = 1}^{\infty} \binom{\alpha}{n} x^n \\ 
                             & = f(x)
\end{align*}
In other words, $(1 + x)f'(x) = \alpha f(x)$. 
Now, using the differential relation,  
\begin{align*}
  \frac{d}{dx}\left[ \frac{f(x)}{(1 + x)^\alpha} \right] 
  & = \frac{f'(x) (1 + x)^\alpha - \alpha f(x) (1 + x)^{\alpha - 1}}{(1 + x)^{2\alpha}} \\ 
  & = \frac{f'(x) (1 + x) - \alpha f(x)}{(1 + x)^{\alpha + 1}} \\ 
  & = \frac{\alpha f(x) - \alpha f(x)}{(1 + x)^{\alpha + 1}} \\ 
  & = 0 
\end{align*}
So $f(x) = c(1 + x)^\alpha$. 
Plugging in $x = 0$ gives $c = 1$, which gives the desired result. 

\end{document}
